\begin{textsample}{NEG Dim 1 – human – Score -7.00 – sp/sp\_ef\_linguagens\_ed\_fis\_1\_p1.txt}  \label{ex:f1_neg_006}
Amparado pela perspectiva cultural, o ensino de Educação Física busca a compreensão do sujeito inserido em diferentes realidades culturais nas quais corpo, \textbf{movimento} e intencionalidade são indissociáveis, o que sugere, para além da vivência, a valorização e a fruição das práticas \textbf{corporais}, bem como a identificação dos \textbf{sentidos} e \textbf{significados} \textbf{produzidos} por estas nos diversos contextos.
Nessa perspectiva, portanto, o currículo deve refletir o contexto sócio histórico : a instabilidade da dinâmica social contemporânea imprime a necessidade de rever, ressignificar e atualizar a visão de cidadão que se pretende formar, bem como os conhecimentos, métodos e o tipo de organização escolar que correspondem a essa formação.

% matched lemmas: corporal, movimento, produzir, sentido, significado
\end{textsample}
